\documentclass[aps,prd,twocolumn,showpacs,superscriptaddress,nofootinbib]{revtex4-2}

\usepackage{amsmath,amssymb}
\usepackage{graphicx}
\usepackage{booktabs}
\usepackage{siunitx}
\usepackage{hyperref}

% Canonical counts pipeline (Phase 5v: freshness gate).
% Auto-generated by update_counts.py — 2026-04-24T19:44:03.140400
% Do not edit manually. Run: uv run python scripts/update_counts.py
\newcommand{\totaltheorems}{3993}
\newcommand{\substantivetheorems}{3882}
\newcommand{\placeholdertheorems}{111}
\newcommand{\axiomcount}{1}
\newcommand{\sorrycount}{0}
\newcommand{\sorrypercent}{0.0}
\newcommand{\leanmodules}{167}
\newcommand{\leandefinitions}{3297}
\newcommand{\pythonmodules}{68}
\newcommand{\testfiles}{59}
\newcommand{\totaltests}{2836}
\newcommand{\figurecount}{111}
\newcommand{\notebookcount}{54}
\newcommand{\papercount}{19}
\newcommand{\aristotleproved}{322}
\newcommand{\aristotleruns}{44}
% --- Paper 18 / Phase 5t per-section counts (FermiHubbardDimer.lean) ---
\newcommand{\fhdTotal}{143}
\newcommand{\fhdWaveTwo}{13}
\newcommand{\fhdWaveThree}{6}
\newcommand{\fhdWaveFour}{22}
\newcommand{\fhdWaveFive}{22}
\newcommand{\fhdWaveSixABC}{38}
\newcommand{\fhdWaveSixExtra}{15}
\newcommand{\fhdWaveSeven}{19}
\newcommand{\fhdWaveEight}{8}


\newcommand{\HiggsField}{\Phi}
\newcommand{\TempCrit}{T_c}
\newcommand{\CubicCoeff}{E}

\begin{document}

\title{Electroweak Phase Transition Order from a Microscopic Scalar
Rung:\\
A Lean-formalized prediction with EW baryogenesis viability}

\author{John G.\ Roehm}
\affiliation{NetRxn Foundation}

\date{\today}

\begin{abstract}
The order of the electroweak phase transition (first-order vs.\ crossover)
is a microscopic prediction of the scalar-rung condensate hierarchy on
the WetterichNJL substrate. Building on the Wave 1b BHL gauge-embedding
formalization (\texttt{BHLGaugeEmbedding.lean}, $23$ theorems) — itself
an extension of the Bardeen-Hill-Lindner / Hill 2025 bilocal mechanism
onto the Clifford-16 Fierz basis~\cite{BardeenHillLindner,Hill2025Redux}
— we ship \texttt{EWPhaseTransition.lean}, $21$ theorems formalizing
the leading-order high-T effective potential, the first-order /
crossover predicate, the latent-heat order parameter, and the EW
baryogenesis viability marker. The transition order is determined
microscopically by the cubic coefficient $E$: first-order if $E > 0$
(barrier survives), crossover if $E = 0$.

\textit{At the strict-leading-order high-T expansion}, the SM benchmark
($m_H = 125.20$ GeV, $\lambda_{\rm SM} \approx 0.13$, $E_{\rm SM} \approx
0.01$) yields $E > 0$ and is therefore predicted first-order; this is
the load-bearing Lean theorem \texttt{sm\_benchmark\_is\_first\_order}.
The physical SM, however, is a \textit{crossover}: lattice analyses
demonstrate that thermal-loop corrections to $E$ at the physical
$m_H = 125.20$ GeV~\cite{PDG2024} push the effective cubic coefficient
below the threshold needed for a first-order transition. The
crossover-to-first-order endpoint sits at $m_H \approx 72.4 \pm 1.7$
GeV, refined by the lattice study of Csikor, Fodor, and
Heitger~\cite{CsikorFodorHeitger1999} (extending the earlier $70$--$95$
GeV range from~\cite{KLRS1996}); the physical $m_H$ is well above
this value. The strict-LO Lean prediction (first-order) and the
physical lattice verdict (crossover) thus diverge at a known and
literature-documented boundary. The Lean theorem
\texttt{crossover\_excludes\_baryogenesis} formalizes the consequence:
a crossover transition is not viable for sphaleron-decoupling
baryogenesis, forcing baryogenesis to leptogenesis or BSM mechanisms ---
the load-bearing input to the future Phase 6c.2 EW-baryogenesis bridge
theorem. Two tracked hypotheses
(\texttt{H\_VacuumStableUnderRG}, \texttt{H\_HierarchyEWLambdaUV})
encode the SM meta-stability and EW-Λ_UV hierarchy as falsifiable
parametric conditions; both have explicit non-trivial witnesses and
falsifiers. The full formalization is zero-sorry, introduces zero new
axioms, and adds zero heartbeat overrides.
\end{abstract}

\pacs{12.15.-y, 12.60.Cn, 11.10.Wx, 11.30.Qc}

\maketitle

%% ================================================================
\section{Introduction}
\label{sec:intro}

The order of the electroweak phase transition is a quantitative
prediction of the SM Higgs sector at finite temperature. The
Kajantie-Laine-Rummukainen-Shaposhnikov 1996 lattice
analysis~\cite{KLRS1996} established that, \emph{after thermal-loop
resummation}, the SM is a crossover for $m_H \gtrsim 70$--$95$ GeV,
with the endpoint refined by Csikor, Fodor, and
Heitger~\cite{CsikorFodorHeitger1999} to $m_H = 72.4 \pm 1.7$
GeV; with the PDG observed $m_H = 125.20$
GeV~\cite{PDG2024}, the resummed SM transition is a crossover, and EW
baryogenesis in the strict-SM (without BSM extensions of the
electroweak Higgs sector) is excluded. The strict-LO Lean prediction
in this paper instead encodes the un-resummed cubic coefficient and
is reframed below as a structural reference point against which the
resummed lattice result improves.

The Phase 5z Wave 3 program formalizes this prediction at the level
of the WetterichNJL scalar-rung embedding (Wave 1, with BHL gauge
extension Wave 1b). The leading-order high-T expansion of the Higgs
effective potential takes the form
\begin{equation}
\label{eq:vt-expansion}
V_T(\HiggsField, T) = \tfrac{1}{2} (c_T T^2 - \mu^2) \HiggsField^2
                       - \tfrac{1}{3} \CubicCoeff T \HiggsField^3
                       + \tfrac{1}{4} \lambda \HiggsField^4,
\end{equation}
where $c_T$ is the thermal-mass coefficient ($\sim (g^2 + g'^2 + 4 y_t^2)/16$
in the SM), $\CubicCoeff$ is the cubic coefficient from
hard-thermal-loop W/Z resummation ($\sim g^3$), and the cubic term is
what distinguishes a first-order transition (with a barrier between
false and true vacua) from a crossover (smooth interpolation).

This paper formalizes \eqref{eq:vt-expansion} as
\texttt{EWPhaseTransition.finiteTPotential}, the cubic-coefficient
order parameter as \texttt{IsFirstOrderEW} / \texttt{IsCrossoverEW},
the latent-heat consequence as \texttt{latentHeat\_nonneg} +
\texttt{crossover\_has\_zero\_latent\_heat} +
\texttt{first\_order\_has\_positive\_latent\_heat}, and the
baryogenesis viability marker as
\texttt{IsBaryogenesisViable} +
\texttt{crossover\_excludes\_baryogenesis} --- in a single Lean module
of $\substantivetheorems{}$-substantive-theorem project context.

%% ================================================================
\section{Finite-T Effective Potential and Order Parameter}
\label{sec:finite-T-potential}

\subsection{Structure invariants}

The Lean structure \texttt{EWFiniteTParams} carries the four LO-expansion
parameters $(\mu^2, \lambda, c_T, \CubicCoeff)$ together with positivity
or non-negativity invariants on each:
$\mu^2 > 0$, $\lambda > 0$, $c_T > 0$, $\CubicCoeff \geq 0$. The
non-negativity (rather than strict positivity) of $\CubicCoeff$ is
load-bearing for the order classifier: it allows the partition into
\textit{first-order} ($\CubicCoeff > 0$) and \textit{crossover}
($\CubicCoeff = 0$) without permitting the unphysical $\CubicCoeff < 0$
branch (which the high-T expansion does not produce at the SM gauge
content). All downstream substantive theorems (latent-heat positivity,
crossover ↔ zero latent heat) consume this invariant explicitly.

\subsection{High-T expansion}

The thermal mass-squared $\mu^2(T) = c_T T^2 - \mu^2$ defines the
critical temperature
\begin{equation}
T_c = \sqrt{\mu^2 / c_T}
\end{equation}
where $\mu^2(T_c) = 0$. The Lean theorem
\texttt{thermalMassSq\_at\_T\_c} certifies this exactly; below $T_c$,
$\mu^2(T) < 0$ is the broken phase
(\texttt{thermalMassSq\_neg\_below\_T\_c}). At the SM benchmark
($\mu^2 = (88\,{\rm GeV})^2$, $c_T = 0.4$), $T_c \approx 139.13$ GeV,
consistent with literature.

\subsection{Transition-order classifier}

We formalize the LO classifier as the cubic-coefficient sign:
\begin{align}
\texttt{IsFirstOrderEW}(p) &\iff p.\CubicCoeff > 0,\\
\texttt{IsCrossoverEW}(p) &\iff p.\CubicCoeff = 0.
\end{align}
The Lean disjointness theorem
\texttt{first\_order\_and\_crossover\_disjoint} certifies the partition.
The definitional reduction
\texttt{transition\_order\_unfolds\_to\_cubic\_sign} unfolds the order
predicates back to the cubic-coefficient inequality (a purely
syntactic step). The substantive correctness-push content is carried
by two non-trivial biconditionals connecting the microscopic order
parameter to the macroscopic latent-heat observable:
\texttt{first\_order\_iff\_positive\_latent\_heat} (proved via
positivity + non-negativity \texttt{linarith}/\texttt{positivity}
reasoning) and \texttt{latentHeat\_zero\_iff\_crossover}, both of which
require the structural non-negativity invariant
\texttt{cubic\_coeff\_nonneg} of \texttt{EWFiniteTParams}.

\subsection{Latent-heat consequence}

The leading-order latent heat is
\begin{equation}
L = \frac{\CubicCoeff^2 T_c^2}{2 \lambda},
\end{equation}
formalized as \texttt{latentHeat}, which is non-negative
(\texttt{latentHeat\_nonneg}); zero exactly for crossover
(\texttt{crossover\_has\_zero\_latent\_heat}); strictly positive for
first-order with positive $T_c$
(\texttt{first\_order\_has\_positive\_latent\_heat}). The latent heat
is therefore an *order parameter* for the phase transition order: its
vanishing flags the crossover branch.

%% ================================================================
\section{SM Benchmark and Lattice Crossover Threshold}
\label{sec:sm-benchmark}

The Lean structure \texttt{smBenchmarkParams} records the LO SM
benchmark:
\begin{center}
\begin{tabular}{lr}
\toprule
$\mu^2$ & $(88\,\mathrm{GeV})^2 = 7744\,\mathrm{GeV}^2$ \\
$\lambda$ & $0.13$ (PDG 2024) \\
$c_T$ & $0.4$ \\
$\CubicCoeff_{\rm SM}$ & $\approx 0.01$ \\
$T_c$ & $\approx 139.13$ GeV \\
\bottomrule
\end{tabular}
\end{center}

The theorem \texttt{sm\_benchmark\_is\_first\_order} certifies that
\textit{at the strict-LO} thermal expansion, $\CubicCoeff_{\rm SM} > 0$
yields a first-order transition. \textbf{This LO prediction does not
survive resummation:} Kajantie, Laine, Rummukainen, and
Shaposhnikov~\cite{KLRS1996} established via 3D effective-theory lattice
Monte Carlo that the SM EW transition is a crossover for
$m_H \gtrsim 70$--$95$ GeV; Csikor, Fodor, and
Heitger~\cite{CsikorFodorHeitger1999} subsequently refined the
endpoint to $m_H = 72.4 \pm 1.7$ GeV by closing the gap with extended
lattice statistics. With the PDG-observed
$m_H = 125.20 \pm 0.11$~GeV~\cite{PDG2024} well above this endpoint,
the resummed SM is a crossover; the strict-LO prediction encoded in
this paper diverges from the resummed lattice verdict at the
documented threshold. The paper's headline figure marks the
\citeauthor{CsikorFodorHeitger1999} endpoint at $m_H = 72$ GeV, with
the SM star at $m_H = 125.20$ GeV firmly in the crossover region.

\begin{figure}[htbp]
\centering
\includegraphics[width=\columnwidth]{figures/fig_ew_transition_phase_diagram.png}
\caption{\textbf{EW phase transition order classifier.} Left panel:
$V_T(\Phi, T)$ at three temperatures (broken phase $T = 0$, halfway to
critical $T = T_c/2 \approx 70$ GeV, symmetric phase $T = 1.5 T_c$) at
the SM benchmark; curves normalized to unit max-abs for clarity. Right
panel: first-order vs.\ crossover partition over the $(m_H, E)$ plane.
Shaded region ($E > 0$) is first-order at LO; the SM benchmark
(star, $m_H = 125.20$ GeV, $\CubicCoeff = 0.01$) is in the first-order
region at LO, but the Kajantie-Laine-Rummukainen-Shaposhnikov 1996
lattice crossover threshold (dashed cross at $m_H = 72$ GeV) places
the physical SM in the crossover branch once thermal-loop corrections
are resummed. Lean: \texttt{thermalMassSq\_neg\_below\_T\_c},
\texttt{first\_order\_and\_crossover\_disjoint},
\texttt{crossover\_excludes\_baryogenesis}.}
\label{fig:ew_phase}
\end{figure}

%% ================================================================
\section{EW Baryogenesis Viability}
\label{sec:baryogenesis}

EW baryogenesis (Cohen-Kaplan-Nelson 1991) requires a
\textit{strong-enough} first-order transition that
sphaleron processes are decoupled in the broken phase as bubble walls
pass; the canonical threshold is $v(T_c) / T_c > 1$, equivalently
$\CubicCoeff / (\lambda T_c) > 1$ at LO. The predicate
\texttt{IsBaryogenesisViable} formalizes this.

The load-bearing exclusion theorem
\texttt{crossover\_excludes\_baryogenesis} certifies that any
crossover transition fails the baryogenesis viability check, regardless
of the threshold value. In Lean:
\begin{verbatim}
theorem crossover_excludes_baryogenesis
  (p : EWFiniteTParams) (threshold : ℝ)
  (h_co : IsCrossoverEW p) :
    ¬ IsBaryogenesisViable p threshold
\end{verbatim}
The proof discharges by contradiction: a baryogenesis-viable transition
implies first-order; first-order and crossover are disjoint; therefore
no crossover transition can be baryogenesis-viable.

This is the load-bearing input to the future Phase 6c.2 ``EW
baryogenesis $\leftrightarrow$ chirality wall'' bridge theorem: under
the SM-as-crossover verdict, EW baryogenesis is forbidden in the
scalar-rung-embedded SM, pushing baryogenesis to leptogenesis or BSM
mechanisms (consistent with the Phase 5z Wave 2 \texttt{MajoranaRung.lean}
sterile-neutrino seesaw infrastructure).

%% ================================================================
\section{Vacuum Stability and Hierarchy Tracked Hypotheses}
\label{sec:tracked-hypotheses}

Two parametric tracked hypotheses encode SM phenomenology beyond the
LO finite-T expansion:

\textbf{Vacuum stability under RG running} — the SM Higgs quartic
coupling $\lambda$ is known to dip near zero at $\sim 10^{11}$ GeV
under SM RG flow~\cite{ButtazzoEtAl2013}, placing the SM in the
\textit{meta-stable} regime. We encode this as
\texttt{H\_VacuumStableUnderRG}:
$\lambda_{\rm min} \le p.\lambda \land 0 < \lambda_{\rm min}$. The
non-trivial content is that $\lambda_{\rm min}$ is positive — a
falsifying configuration with $\lambda_{\rm min} \le 0$ is rejected by
\texttt{vacuum\_stability\_fails\_at\_negative\_lam\_min}. The SM
benchmark satisfies the hypothesis at $\lambda_{\rm min} = 0.05$
(\texttt{sm\_benchmark\_vacuum\_stable}).

\textbf{EW-$\Lambda_{\rm UV}$ hierarchy} — encoded as
\texttt{H\_HierarchyEWLambdaUV}:
$10^6 v_{\rm EW} < \Lambda_{\rm UV} \land 0 < v_{\rm EW}$. Witness at
GUT scale (\texttt{hierarchy\_holds\_at\_GUT\_scale}); falsifier when
$\Lambda_{\rm UV} \le 10^6 v_{\rm EW}$
(\texttt{hierarchy\_fails\_when\_v\_EW\_too\_large}).

The bundled Wave 3 open manifest
\texttt{Wave3OpenManifest}: vacuum-stability $\land$ hierarchy $\land$
(first-order $\lor$ crossover); witnessed by
\texttt{wave3\_open\_manifest\_consistent} at the SM benchmark with
$\lambda_{\rm min} = 0.05$, $\Lambda_{\rm UV} = 10^{16}$ GeV,
$v_{\rm EW} = 246$ GeV.

%% ================================================================
\section{Bridge to Wave 1b BHL Embedding}
\label{sec:wave1b-bridge}

Wave 3 builds directly on Wave 1b's BHL gauge-embedding transplant:
the doublet Higgs field $H^i = \bar\psi_R \psi_L^i$ identified by the
BHL extension carries SU(2)$_L \times $ U(1)$_Y$ quantum numbers, so the
finite-T potential \eqref{eq:vt-expansion} is well-defined as a
function of $|H|^2 = H^\dagger H$.

The Lean theorem \texttt{ew\_finite\_T\_built\_on\_bhl} certifies the
trivial-zero compatibility (the potential vanishes at $\Phi = 0$ for
any temperature), and \texttt{zero\_temperature\_recovers\_zero\_T\_potential}
certifies that at $T = 0$ the finite-T potential reduces to the
Mexican hat
$V(\Phi) = -\tfrac{1}{2} \mu^2 \Phi^2 + \tfrac{1}{4} \lambda \Phi^4$
of the Wave 1 \texttt{ScalarRungInterpretation.lean} module. This
chain establishes that Wave 3 is a strict refinement of the
scalar-rung embedding: the BHL daughter Higgs doublet of Wave 1b carries
the LO thermal corrections without abandoning the substrate
identification.

%% ================================================================
\section{Formal Verification Summary}
\label{sec:formal-verification}

The Wave 3 formalization comprises one new Lean module:
\begin{center}
\begin{tabular}{ll}
\toprule
Module & Theorems \\
\midrule
\texttt{EWPhaseTransition.lean} & 21 \\
\bottomrule
\end{tabular}
\end{center}
plus one structure (\texttt{EWFiniteTParams}), $12$ non-computable
definitions for the potential / mass / $T_c$ / latent-heat / order-parameter
classifiers, and four tracked-hypothesis Props
(\texttt{H\_VacuumStableUnderRG}, \texttt{H\_HierarchyEWLambdaUV},
\texttt{IsBaryogenesisViable}, \texttt{Wave3OpenManifest}). The module
is zero-sorry, introduces zero new axioms, and adds zero heartbeat
overrides (Phase~3 pipeline invariant 10).

The module integrates with three prior-phase modules:
\texttt{ScalarRungInterpretation.lean} (Wave 1; Mexican-hat at $T = 0$),
\texttt{BHLGaugeEmbedding.lean} (Wave 1b; doublet Higgs identification),
and \texttt{WetterichNJL.lean} (substrate Fierz infrastructure).
A Python companion package \texttt{src/ew\_phase\_transition/} provides
numerical wrappers for the finite-T potential, transition-order
classifier, and baryogenesis-viability marker; canonical formulas
(\texttt{ew\_finite\_t\_potential}, \texttt{ew\_thermal\_mass\_sq},
\texttt{ew\_critical\_temperature}, \texttt{ew\_latent\_heat}) live in
\texttt{src/core/formulas.py} per pipeline invariant. A 36-test pytest
suite covers every Lean theorem; full-project pytest passes.

\section{Outlook}
\label{sec:outlook}

Phase 6c.2 (``EW baryogenesis $\leftrightarrow$ chirality wall bridge'')
inherits the
\texttt{crossover\_excludes\_baryogenesis} verdict as a hard
constraint: in the SM-as-embedded-in-scalar-rung framing, EW
baryogenesis is forbidden, motivating the leptogenesis branch via
the Phase 5z Wave 2 \texttt{MajoranaRung.lean} sterile-neutrino seesaw.
Two-loop thermal corrections to $\CubicCoeff$ (deferred from this
wave's strict-LO scope) and Boltzmann-code nucleation rates (deferred
to Phase 6b.3) close the quantitative scope of the SM transition
order. The vacuum-stability tracked hypothesis is the natural target
for an asymptotic-safety closure (Shaposhnikov-Wetterich
2010~\cite{ShaposhnikovWetterich2010}), formalizing the IR fixed-point
prediction $m_H = 126 \pm $ few GeV from a different (elementary-Higgs)
mechanism.

%% ================================================================
\begin{thebibliography}{99}

\bibitem{BardeenHillLindner}
W.~A.~Bardeen, C.~T.~Hill, and M.~Lindner,
\emph{Minimal dynamical symmetry breaking of the standard model},
Phys.~Rev.~D~\textbf{41}, 1647 (1990); DOI:~\href{https://doi.org/10.1103/PhysRevD.41.1647}{10.1103/PhysRevD.41.1647}.

\bibitem{Hill2025Redux}
C.~T.~Hill,
\emph{Natural top quark condensation (a redux)},
arXiv:2503.21518 (2025); \url{https://arxiv.org/abs/2503.21518}.

\bibitem{KLRS1996}
K.~Kajantie, M.~Laine, K.~Rummukainen, and M.~Shaposhnikov,
\emph{Is there a hot electroweak phase transition at $m_H \gtrsim m_W$?},
Phys.~Rev.~Lett.~\textbf{77}, 2887 (1996); arXiv:hep-ph/9605288;
DOI:~\href{https://doi.org/10.1103/PhysRevLett.77.2887}{10.1103/PhysRevLett.77.2887}.

\bibitem{CsikorFodorHeitger1999}
F.~Csikor, Z.~Fodor, and J.~Heitger,
\emph{Endpoint of the hot electroweak phase transition},
Phys.~Rev.~Lett.~\textbf{82}, 21 (1999); arXiv:hep-ph/9809291;
DOI:~\href{https://doi.org/10.1103/PhysRevLett.82.21}{10.1103/PhysRevLett.82.21}.

\bibitem{ButtazzoEtAl2013}
D.~Buttazzo, G.~Degrassi, P.~P.~Giardino \textit{et al.},
\emph{Investigating the near-criticality of the Higgs boson},
JHEP \textbf{12} (2013) 089; arXiv:1307.3536.

\bibitem{ShaposhnikovWetterich2010}
M.~Shaposhnikov and C.~Wetterich,
\emph{Asymptotic safety of gravity and the Higgs boson mass},
Phys.~Lett.~B~\textbf{683}, 196 (2010); arXiv:0912.0208.

\bibitem{PDG2024}
S.~Navas \textit{et al.}~(Particle Data Group),
\emph{Review of particle physics},
Phys.~Rev.~D~\textbf{110}, 030001 (2024).

\bibitem{Lean4}
L.~de~Moura and S.~Ullrich,
\emph{The Lean 4 theorem prover and programming language},
in \emph{Automated Deduction---CADE 28} (Springer, 2021), pp.~625--635.

\bibitem{Mathlib}
The Mathlib Community,
\emph{The Lean mathematical library},
in \emph{Proceedings of the 9th ACM SIGPLAN International Conference on Certified Programs and Proofs} (CPP 2020).

\end{thebibliography}

\end{document}
